\scanpage{24}
\mspage{- 2 -}
\noindent\hfill n\textsuperscript{o}355

d'o\handwritten{\`}{u} $\operatorname{Ass}_B(F\handwritten{/}\underline p F)\handwritten{\subset}\operatorname{Ass}_B(E\handwritten{\otimes}_A F)$, ce qui démontre l'inclusion $(\ast)$.

On va maintenant supposer que $A$ est \uline{noethérien}, et démontrer l'inclusion opposée. On procèdera par étapes :

(i) $E$ \uline{est un module de type fini et} $\operatorname{Ass}_A(E)$ \uline{est réduit à un élément} $\underline p$.

D'après le th. 1 du \handwritten{\S}{} 1, il existe une suite de composition $(E_i)_{0\handwritten{\leq} i\handwritten{\leq} n}$ de $E$ telle que $E_i\handwritten{/}E_{i+1}$ soit isomorphe à $A\handwritten{/}\underline p_i$, o\handwritten{\`}{u} $\underline p_i$ est un idéal premier de $A$. D'après le th.2 du \handwritten{\S}{} 1, tous les $\underline p_i$ contiennent $\underline p$. Du fait que $F$ est $A$-plat, les $E_i\handwritten{\otimes}_A F$ forment une suite de composition de $E\handwritten{\otimes}_A F$, et $(E_i\handwritten{\otimes}_A F)\handwritten{/}(E_{i+1}\handwritten{\otimes}_A F)$ s'identifie à $A\handwritten{/}\underline p_i\handwritten{\otimes}_A F=F\handwritten{/}\underline p_i F$. En appliquant la prop.3, on en déduit :

\[
\operatorname{Ass}_B(E\handwritten{\otimes}_A F)\handwritten{\subset}\handwritten{\bigcup}\operatorname{Ass}_B(F\handwritten{/}\underline p_i F)\,.
\]

Soit donc $\underline P\handwritten{\in}\operatorname{Ass}_B(E\handwritten{\otimes}_A F)\handwritten{\cap}\operatorname{Ass}_B(F\handwritten{/}\underline p_i F)$. Si $\underline p_i\handwritten{\ne}\underline p$, il existe $a\handwritten{\in}\underline p_i$, $a\handwritten{\notin}\underline p$. D'après le cor.2 à la prop.2 du \handwritten{\S}{} 1, l'homothétie de rapport $a$ est injective dans $E$, donc celle de rapport $\handwritten{\pi}(a)$ est injective dans $E\handwritten{\otimes}_A F$ ; mais d'autre part, $\handwritten{\pi}(a)$ annule $F\handwritten{/}\underline p_i F$, d'o\handwritten{\`}{u} $\handwritten{\pi}(a)\handwritten{\in}\underline P$, contradiction. Ainsi $\underline p_i$ est nécessairement égal à $\underline p$, ce qui démontre la proposition dans le cas considéré.

(ii) $E$ \uline{est un module de type fini}.

Soient $\underline p_i$ les éléments de $\operatorname{Ass}_A(E)$, et soit $0=\handwritten{\bigcap} Q_i$ une décomposition primaire réduite correspondante (cf. \handwritten{\S}{} 2, n\textsuperscript{o}2). Le module $E$ se plonge dans la somme directe des $E_i=E\handwritten{/}Q_i$, et $E\handwritten{\otimes}_A F$ se plonge dans la somme directe des $E_i\handwritten{\otimes}_A F$. On en déduit :

\[
\operatorname{Ass}_B(E\handwritten{\otimes}_A F)\handwritten{\subset}\handwritten{\bigcup}\operatorname{Ass}_B(E_i\handwritten{\otimes}_A F)\,.
\]

D'après (i), on a $\operatorname{Ass}_B(E_i\handwritten{\otimes}_A F)=\operatorname{Ass}_B(F\handwritten{/}\underline p_i F)$, d'o\handwritten{\`}{u} la proposition dans ce cas.

(iii) \uline{Cas général}.

% archive sheet number: 23
